\documentclass[11pt,a4paper]{article}
\usepackage{curso}
\begin{document}
\portada{Lección}{0.1}{El muro energético de la IA}

\noindent\textbf{Unidad 0 --- Orientación y nivelación matemática.}\quad
\textbf{Objetivos:} (i) derivar la energía de conmutación $\tfrac12CV^2$ y la potencia dinámica
$P=\alpha CV^2f$; (ii) explicar el escalado de Dennard y su final; (iii) argumentar por qué mover datos
cuesta más que calcular; (iv) usar el modelo roofline.
\textbf{Prerrequisitos:} circuitos básicos (capacitor, resistencia), notación científica.

\section{¿Por qué importa?}
La pregunta central del curso es cuándo, por qué y en cuánto el hardware probabilístico basado en la
física puede hacer la IA más eficiente energéticamente. Para responderla primero hay que saber
\emph{dónde se gasta hoy la energía}. El pensum lo resume así: la capacidad de la red eléctrica es ya
un obstáculo principal para construir centros de datos, y la mayor parte de la energía de un acelerador
digital se gasta en \emph{mover} bits y en \emph{suprimir} el ruido térmico, no en calcular \cite{conte}.

\begin{intuicion}
Una fábrica eficiente con el almacén al otro lado de la ciudad: producir es barato, el transporte es
caro. En un chip el núcleo aritmético es la fábrica, la memoria es el almacén y los cables son la
carretera. Cada cable es un capacitor que hay que cargar y descargar.
\end{intuicion}

\section{Energía de conmutación}
Un nodo digital (la salida de una compuerta más el cable que alimenta) se modela como un capacitor $C$
que una fuente de voltaje $V$ carga a través de la resistencia $R$ del transistor.
\begin{proposicion}
Al cargar $C$ de $0$ a $V$ desde una fuente fija, la fuente entrega $E_{\text{fuente}}=CV^2$, el capacitor
almacena $\tfrac12CV^2$ y la resistencia disipa $\tfrac12CV^2$, \emph{independientemente de $R$}.
\end{proposicion}
\begin{proof}
La carga transferida es $Q=CV$ a potencial constante $V$: $E_{\text{fuente}}=QV=CV^2$. La energía en el
capacitor es $\int_0^{Q}\frac{q}{C}\dd q=\tfrac12CV^2$. Directamente: $i(t)=\frac VR e^{-t/RC}$, y
$E_R=\int_0^\infty Ri^2\dd t=\frac{V^2}{R}\int_0^\infty e^{-2t/RC}\dd t=\frac{V^2}{R}\cdot\frac{RC}{2}=\tfrac12CV^2$.
\end{proof}
Al descargar (transición $1\to0$) la energía almacenada también se disipa, de modo que un ciclo
completo $0\to1\to0$ cuesta $CV^2$. Si en promedio una fracción $\alpha$ (factor de actividad) de la
capacidad total conmutada $C$ completa un ciclo cada periodo de reloj,
\begin{resultado}[ (potencia dinámica)]
\begin{equation}P_{\text{din}}=\alpha\,C\,V^2 f .\end{equation}
\end{resultado}
\noindent(Algunos textos definen $\alpha$ por transición y escriben $\tfrac12\alpha CV^2f$; la física es la misma.)
A esto se suma la potencia estática por fugas, que crece al bajar el voltaje umbral.

\section{Escalado de Dennard y su final}
Dennard et al. (1974) observaron que si todas las dimensiones y voltajes de un MOSFET se reducen por un
factor $k>1$, entonces $C\to C/k$, $V\to V/k$, el retardo baja por $k$ y $f\to kf$. Así:
\[
P_{\text{transistor}}\propto CV^2f\to \frac{C}{k}\frac{V^2}{k^2}kf=\frac{P}{k^2},\qquad
\frac{\text{transistores}}{\text{área}}\to k^2\cdot\frac{\text{transistores}}{\text{área}},\qquad
\frac{P}{\text{área}}\to\text{constante}.
\]
Hacia mediados de la década de 2000 el voltaje de alimentación dejó de bajar: reducir el voltaje umbral
dispara la corriente de fuga (que crece exponencialmente al bajar el umbral) y estrecha los márgenes de
ruido. Con $V$ fijo, $P_{\text{transistor}}\to(C/k)V^2(kf)=P$ y la densidad de potencia crece como $k^2$
por generación. Consecuencias: frecuencias estancadas, «silicio oscuro» (no se pueden encender todos los
transistores a la vez) y la potencia como recurso escaso. La ley de Moore (densidad) continuó; el
escalado de Dennard (potencia por área constante) no.

\section{El costo de mover datos}
La capacidad de un cable crece con su longitud. Con un valor típico de orden de magnitud \evE{} de
$C'\approx\SI{0.2}{fF/\micro m}$, un cable de 1 mm tiene $C\approx\SI{200}{fF}$ y a $V=0{,}8$ V una
transición disipa $\tfrac12CV^2\approx\SI{64}{fJ}$, unas 200 veces más que una compuerta local de
$\sim1$ fF. Las mediciones clásicas de Horowitz para 45 nm (referencia complementaria, fuera del pensum)
dan del orden de 0,9 pJ para una suma de punto flotante de 32 bits, 3,7 pJ para una multiplicación,
5 pJ para leer 32 bits de una SRAM de 8 KB y 640 pJ para leer 32 bits de DRAM externa.
\begin{ejemplo}
Un acelerador ejecuta un producto matriz--vector $y=Wx$ con $W\in\R^{4096\times4096}$ en fp16, en lote 1.
Operaciones: $2\cdot4096^2\approx3{,}4\times10^7$ FLOP. Bytes leídos: $2\cdot4096^2\approx3{,}4\times10^7$ B.
Intensidad aritmética $I\approx1$ FLOP/byte. Si leer un byte de DRAM cuesta del orden de 160 pJ ((640 pJ / 32 bits) $\times$ 8) y una operación
$\sim1$ pJ, la energía de memoria domina por un factor del orden de $10^2$.
\end{ejemplo}

\section{El cuello de botella de von Neumann y el roofline}
En la arquitectura de von Neumann memoria y procesador están separados. El modelo roofline acota el
rendimiento alcanzable:
\begin{equation}\text{rend.}(I)=\min\bigl(\pi_{\max},\ \beta_{\text{mem}}\,I\bigr),\qquad I^*=\pi_{\max}/\beta_{\text{mem}},\end{equation}
donde $\pi_{\max}$ es el pico de cómputo (FLOP/s), $\beta_{\text{mem}}$ el ancho de banda (B/s) e $I$ la
intensidad aritmética (FLOP/B). Para $I<I^*$ el programa está limitado por memoria. La decodificación
autorregresiva de un LLM con lote 1 tiene $I\approx1$ (cada peso se usa una vez) y queda muy a la
izquierda del codo; volveremos a esto en la Unidad 10 (prefill frente a decode).

\begin{aplicacion}
\textbf{¿A dónde va la energía de una GPU?} (i) Movimiento de datos entre registros, SRAM, HBM y red;
(ii) márgenes de voltaje para que cada bit sea determinista frente al ruido; (iii) generación de
aleatoriedad pseudoaleatoria para cargas probabilísticas (difusión, muestreo de tokens, inferencia
bayesiana). El reto del laboratorio L1 mide con NVML la energía por muestra aleatoria (J/muestra), que
será la línea base digital de todo el semestre.
\end{aplicacion}

\section{Conexiones}
La Unidad 1 mostrará que el límite fundamental por bit borrado es $\kB T\ln2\approx2{,}9\times10^{-21}$ J,
muy por debajo de $CV^2$. La Unidad 5 introduce la contabilidad
$E_{\text{total}}=E_{\text{tejido}}+E_{\text{E/S}}+E_{\text{digital}}$, y la Unidad 10 aplica el roofline a
prefill y decode en transformers.

\begin{error}
\begin{itemize}[nosep]
\item La disipación por transición no depende de $R$: es $\tfrac12CV^2$ (cambia solo cuán rápido ocurre).
\item Contar FLOPs no basta: en cargas con baja intensidad aritmética domina el movimiento de bytes.
\item Moore $\neq$ Dennard: la densidad siguió creciendo; la potencia por área dejó de ser constante.
\item Potencia (W) no es energía (J): compare energía por tarea (J/token, J/muestra), no vatios.
\end{itemize}
\end{error}

\section{Resumen}
\begin{itemize}[nosep]
\item Ciclo de un nodo: $CV^2$; potencia dinámica $P=\alpha CV^2f$.
\item El fin de Dennard convirtió la potencia en el recurso escaso.
\item Mover datos cuesta órdenes de magnitud más que calcular; von Neumann separa ambos.
\end{itemize}

\begin{thebibliography}{9}\small
\bibitem[6]{conte} T. Conte et al., \emph{Thermodynamic Computing} (CCC workshop report), 2019. \url{https://arxiv.org/abs/1911.01968}
\bibitem[7]{quanta} Quanta Magazine, \emph{Thermodynamic Computers Go With the (Energy) Flow}, 2026.
\bibitem[21]{extropic} Extropic, \emph{Thermodynamic Computing From Zero to One}, 2025. \url{https://extropic.ai/writing}
\bibitem[c]{horowitz} (Complementaria, fuera del pensum) M. Horowitz, Computing's energy problem (and what we can do about it), ISSCC 2014.
\bibitem[c]{dennard} (Complementaria, fuera del pensum) R. H. Dennard et al., Design of ion-implanted MOSFETs with very small physical dimensions, \emph{IEEE J. Solid-State Circuits} 9, 256 (1974).
\end{thebibliography}
\end{document}
